\SourceLeaf{151}
von 1794. Hätte man, in Beziehung auf sämmtliche 204 der Witwenkasse bis jetzt beigetretene Professoren regelmässig aufgezeichnet, ob und während welches Theils ihrer Genossenschaft sie verehelicht gewesen sind, zugleich mit der genauen Altersangabe für sie selbst und ihre Frauen, so würde dieses vermittelst der Mortalitätstafeln zu einer sehr genauen indirecten Bestimmung des fraglichen Verhältnisses benutzt werden können. Von allem dem ist aber, meines Wissens, Nichts geschehen. Je mehr diess jetzt beklagt werden muss, desto zuversichtlicher darf wohl gehofft werden, dass durch zweckmässige wenigstens von jetzt an zu treffende Maassregeln, demjenigen, welcher, wieder nach 100 Jahren, begutachten wird, eine gleiche Klage erspart sein wird.

Soviel über das \textsc{Kritter}’sahe Gutachten, oder vielmehr über denjenigen Theil desselben, der mit dem Gegenstande meiner gegenwärtigen Kritik in unmittelbarer Verbindung steht; auf den zweiten Theil des Gutachtens, der die aus den Waisenpensionen entspringende Vergrösserung der Ausgaben betrifft, werde ich in der zweiten Abtheilung dieser Denkschrift zurückkommen. Dass in jenem ersten Theile des \textsc{Kritter}’schen Gutachtens eben nur die Oberfläche des Gegenstandes berührt ist, wird, meine ich, durch die vorstehenden Entwicklungen zur Genüge dargethan sein. Betrachtungen oder Vermuthungen darüber, wie es zugegangen, dass \textsc{Kritter} nur ein so oberflächliches, nicht einmal mit seinen sonstigen öffentlichen Äusserungen übereinstimmendes Gutachten ausgestellt hat, würden mich hier zu weit führen. Indess möchte ich glauben, dass, wenn man, anstatt sich auf die Verlegung von zwei Specialfragen zu beschränken,
\begin{quote}
wovon zudem die eine, in dem Sinn wie der Fragesteller meinte, eine bestimmte Antwort gar nicht zuliess, den p. \textsc{Kritter}, gleichviel in welcher Form, an den weitern Deliberationen hätte Theil nehmen lassen, oder ihn wenigstens über die Angemessenheit des Plans, den man auf sein Gutachten gründen zu können vermeinte, unter vollständiger Mittheilung aller Sachverhältnisse zu Rathe gezogen hätte,
\end{quote}
\noindent die Frage über die Normirung aller künftigen Pensionserhöhungen nicht so, wie geschehen, über das Knie gebrochen sein würde\footnote{Die Zahl der im Laufe meines obigen Berichts angeführten oder angedeuteten Züge von laxer Geschäftsbehandlung hätte leicht noch vergrössert werden können, was ich jedoch für so unnöthig wie unerfreulich gehalten habe.}. Das ist aber unterblieben. \textsc{Kritter} erhielt aus der Witwenkasse ein Honorarium von 4 Thl. 24 mgr. für sein Gutachten, so wie P., in dem Ministerialrescript vom 20. November 1794, eine Belobung seiner wohlausgearbeiteten Denkschrift.

In diesem Rescript wurden die gemachten Vorschläge genehmigt, wegen der Erhöhung der Pensionen jedoch bevorwortet, dass, wenn wider Verhoffen \textit{unglückliche Umstände} demnächst eine Verminderung der Pensionen nothwendig machen sollten, die alsdann vorhandenen Witwen sich solches gefallen lassen müssen; auch wurde für die künftig der Progressionsnormirung gemäss vorzunehmenden Pensionserhöhungen die jedesmalige Ratification vorbehalten. \textsuperscript{.} Es scheint bemerkenswerth, dass hier nur einer solchen Nothwendigkeit gedacht ist, die aus unglücklichen Umständen, nicht aber derjenigen, die möglicherweise aus einer zu grossen Witwenzahl hervorgehen könnte. Hat man einen solchen Fall für unmöglich gehalten, oder hat man das Regulativ in demselben Sinn wie \textsc{Sch.} aufgefasst, und, dass in einem solchen Fall die Pension wieder herabgehen müsse, als sich von selbst verstehend betrachtet? Übrigens wurde bei Ratification der ersten Erhöhung (1799 Mai 23) derselbe Vorbehalt wiederholt, aber nur im Allgemeinen von Umständen, die die Wiederverminderung nothwendig machen könnten, gesprochen, ohne die Qualification von unglücklichen. In den spätern Ratificationsfällen ist, so viel ich habe finden können, die Reservation nicht wiederholt.
